\section{Preliminaries}\label{sec:prelim}

This section fixes the terminology and notation and introduces the running
example.

\subsection{Graphs and paths}\label{sec:def-graph}

\begin{definition}[Weighted graph]\label{def:graph}
A \emph{weighted graph} $G = (V, E, w)$ consists of a finite set $V$ of
vertices, a set $E$ of edges, each of which joins two vertices, and a weight
function $w$ that assigns a nonnegative cost $w(u,v)$ to every edge $(u,v)$.
\end{definition}

All graphs in this report are undirected. In Figure~\ref{fig:defs},
$V = \{A, p, q, r, B\}$ and $w(A,p) = 2$.

\begin{definition}[Path and path cost]\label{def:path}
A \emph{path} from a vertex $u$ to a vertex $v$ is a sequence of vertices
$P = (v_0, v_1, \dots, v_k)$ with $v_0 = u$ and $v_k = v$ in which every two
consecutive vertices are joined by an edge. The \emph{cost} of the path is
\begin{equation}\label{eq:pathcost}
  c(P) = \sum_{i=1}^{k} w(v_{i-1}, v_i).
\end{equation}
\end{definition}

For example, the path $(A, q, r, B)$ of Figure~\ref{fig:defs-b} costs
$4 + 2 + 3 = 9$.

\begin{definition}[Distance and shortest path]\label{def:shortest}
The \emph{distance} $\delta(u,v)$ from a vertex $u$ to a vertex $v$ is the
minimum cost of a path from $u$ to $v$; if no such path exists, then
$\delta(u,v) = \infty$. A path from $u$ to $v$ of cost $\delta(u,v)$ is a
\emph{shortest path} from $u$ to $v$.
\end{definition}

Table~\ref{tab:defs-paths} lists the paths from $A$ to $B$ that visit no vertex
twice; the shortest one has the most edges.

\begin{figure}[htbp]
  \centering
  \newcommand{\defsgraph}[2]{%
  \begin{tikzpicture}[x=0.9cm, y=0.9cm]
    \coordinate (A) at (0,0.9);   \coordinate (p) at (1.6,1.8);
    \coordinate (q) at (1.6,0);   \coordinate (r) at (3.2,0.9);
    \coordinate (B) at (4.8,0.9);
    \ifx\relax#1\relax\else
      \draw[#2, line width=5pt, line cap=round, line join=round] #1;
    \fi
    \foreach \a/\b/\w in {A/p/2, A/q/4, p/q/1, p/r/5, q/r/2, r/B/3}
      {\draw[storyedge] (\a) -- node[storyweight] {\w} (\b);}
    \node[storystart] at (A) {$A$};
    \node[storynode]  at (p) {$p$};
    \node[storynode]  at (q) {$q$};
    \node[storynode]  at (r) {$r$};
    \node[storygoal]  at (B) {$B$};
  \end{tikzpicture}}
\begin{subfigure}[t]{0.32\textwidth}\centering
  \defsgraph{}{}
  \caption{A weighted graph with five vertices and six edges.}\label{fig:defs-a}
\end{subfigure}\hfill
\begin{subfigure}[t]{0.32\textwidth}\centering
  \defsgraph{(A) -- (q) -- (r) -- (B)}{alggreedy!45}
  \caption{The path $(A,q,r,B)$ has cost $4 + 2 + 3 = 9$.}\label{fig:defs-b}
\end{subfigure}\hfill
\begin{subfigure}[t]{0.32\textwidth}\centering
  \defsgraph{(A) -- (p) -- (q) -- (r) -- (B)}{algastar!45}
  \caption{The shortest path $(A,p,q,r,B)$ has cost $2 + 1 + 2 + 3 = 8$.}\label{fig:defs-c}
\end{subfigure}

  \caption{Paths in the road map of Figure~\ref{fig:pf-graph}.}\label{fig:defs}
\end{figure}

\begin{table}[htbp]
  \centering
  \caption{All paths from $A$ to $B$ in Figure~\ref{fig:defs} that visit no
    vertex twice.}\label{tab:defs-paths}
  \begin{tabular}{@{}lcl@{}}
    \toprule
    Path & Edges & Cost\\
    \midrule
    $(A, p, q, r, B)$ & 4 & $2 + 1 + 2 + 3 = \mathbf{8}$ \quad shortest, $\delta(A,B) = 8$\\
    $(A, q, r, B)$    & 3 & $4 + 2 + 3 = 9$\\
    $(A, p, r, B)$    & 3 & $2 + 5 + 3 = 10$\\
    $(A, q, p, r, B)$ & 4 & $4 + 1 + 5 + 3 = 13$\\
    \bottomrule
  \end{tabular}
\end{table}

\subsection{Heuristic functions}\label{sec:def-heuristic}

A search does not know in advance how far each vertex is from the destination
$t$, but it can use a quick estimate, such as the straight-line distance on a
map. We write $h^{*}(v) = \delta(v,t)$ for the true remaining cost.

\begin{definition}[Heuristic function]\label{def:heuristic}
A \emph{heuristic function}, or simply a \emph{heuristic}, is a function $h$
that assigns to every vertex $v$ a nonnegative number $h(v)$, an estimate of
the remaining cost $h^{*}(v)$.
\end{definition}

\subsection{Best-first search}\label{sec:bestfirst}

All three algorithms of this report are \emph{best-first searches}. They keep a
\emph{priority queue}~\cite{baeldungpq} $Q$ of vertices that have been reached but not explored,
each with a numerical \emph{key}; a smaller key means a more promising vertex.
In each round the search removes a vertex $u$ with the smallest key and
\emph{expands} it: it adds $u$ to the set $S$ of \emph{closed} vertices and
examines every edge $(u,v)$, which may insert $v$ into $Q$ or lower its key.
The search stops when it removes the destination. The algorithms differ only in
the key.

The pseudocode uses three queue operations: \textsc{Extract-Min}$(Q)$ removes
and returns a vertex with the smallest key, \textsc{Insert}$(Q, v, k)$ adds $v$
with key $k$, and \textsc{Insert-or-Decrease-Key}$(Q, v, k)$ adds $v$ or, if it
is already in $Q$, lowers its key to $k$. We measure the work of a search by the
number of vertices it expands, the destination included.
Table~\ref{tab:notation} collects the notation.

\begin{table}[htbp]
  \centering
  \caption{Notation used in the report.}\label{tab:notation}
  \begin{tabular}{@{}ll@{}}
    \toprule
    Symbol & Meaning\\
    \midrule
    $G = (V,E,w)$ & weighted graph with vertex set $V$, edge set $E$ and weights $w$\\
    $s$, $t$      & source vertex and destination vertex\\
    $\delta(u,v)$ & distance from $u$ to $v$, the cost of a shortest path\\
    $C^{*}$       & distance from $s$ to $t$, that is, $\delta(s,t)$\\
    $g(v)$        & cost of the cheapest path from $s$ to $v$ found so far\\
    $h(v)$        & heuristic estimate of the remaining cost from $v$ to $t$\\
    $h^{*}(v)$    & true remaining cost $\delta(v,t)$\\
    $f(v)$        & key used by A*, $f(v) = g(v) + h(v)$\\
    $\pi(v)$      & predecessor of $v$ on the cheapest path found so far\\
    $Q$           & priority queue of vertices reached but not yet expanded\\
    $S$           & set of closed vertices, that is, of expanded vertices\\
    \bottomrule
  \end{tabular}
\end{table}

\subsection{The running example}\label{sec:running}

Figure~\ref{fig:running} shows the graph on which we trace the algorithms:
eight cities from Dhaka (\cty{D}) to Sylhet (\cty{S}), with Gazipur (\cty{G}),
Narsingdi (\cty{N}), Kishoreganj (\cty{K}), Brahmanbaria (\cty{B}), Cumilla
(\cty{C}) and Habiganj (\cty{H}).  The edge weights are small integers
chosen for the example, and each disc shows a heuristic value $h(v)$, which never
exceeds the true remaining cost (Table~\ref{tab:story-h}). The shortest path is
$(\cty{D}, \cty{N}, \cty{B}, \cty{H}, \cty{S})$, of cost $C^{*} = 8$.

\begin{figure}[htbp]
  \centering
  \resizebox{0.62\textwidth}{!}{\begin{tikzpicture}[x=1cm, y=1cm]
  \fill[black!10] (-0.529,-0.393) -- (7.052,-1.657) -- (8.132,1.475) -- (0.551,2.739) -- cycle;
  \draw[draw=rblue!14, fill=rbluelt!35, line width=0.6pt] (-0.529,-0.243) -- (7.052,-1.507) -- (8.132,1.625) -- (0.551,2.889) -- cycle;
  \coordinate (D) at (0.432,1.253);
  \coordinate (N) at (2.333,1.598);
  \coordinate (B) at (3.802,0.691);
  \coordinate (H) at (5.702,1.037);
  \coordinate (S) at (7.171,0.130);
  \coordinate (G) at (0.950,0.173);
  \coordinate (K) at (4.234,1.944);
  \coordinate (C) at (5.584,-0.599);
  \draw[draw=rblue!55, line width=1pt] (D) -- (N);
  \draw[draw=rblue!55, line width=1pt] (D) -- (B);
  \draw[draw=rblue!55, line width=1pt] (N) -- (B);
  \draw[draw=rblue!55, line width=1pt] (B) -- (H);
  \draw[draw=rblue!55, line width=1pt] (H) -- (S);
  \draw[draw=rblue!55, line width=1pt] (D) -- (G);
  \draw[draw=rblue!55, line width=1pt] (N) -- (K);
  \draw[draw=rblue!55, line width=1pt] (K) -- (H);
  \draw[draw=rblue!55, line width=1pt] (B) -- (C);
  \draw[draw=rblue!55, line width=1pt] (C) -- (S);
  \node[isoweight] at ($(D)!0.5!(N)$) {2};
  \node[isoweight] at ($(D)!0.5!(B)$) {5};
  \node[isoweight] at ($(N)!0.5!(B)$) {1};
  \node[isoweight] at ($(B)!0.5!(H)$) {2};
  \node[isoweight] at ($(H)!0.5!(S)$) {3};
  \node[isoweight] at ($(D)!0.5!(G)$) {1};
  \node[isoweight] at ($(N)!0.5!(K)$) {3};
  \node[isoweight] at ($(K)!0.5!(H)$) {4};
  \node[isoweight] at ($(B)!0.5!(C)$) {3};
  \node[isoweight] at ($(C)!0.5!(S)$) {10};
  \isodisc{D}{white}{algastar}{D}{7}
  \isodisc{N}{white}{rblue!70}{N}{5}
  \isodisc{B}{white}{rblue!70}{B}{4}
  \isodisc{H}{white}{rblue!70}{H}{2}
  \isodisc{S}{white}{cbRed}{S}{0}
  \isodisc{G}{white}{rblue!70}{G}{8}
  \isodisc{K}{white}{rblue!70}{K}{4}
  \isodisc{C}{white}{rblue!70}{C}{7}
\end{tikzpicture}
}
  \caption{The running example. Each edge shows its weight and each disc the
    heuristic value $h$; the source \cty{D} has a green rim and the destination
    \cty{S} a red rim.}\label{fig:running}
\end{figure}

\begin{table}[htbp]
  \centering
  \caption{The heuristic $h$ of the running example and the true remaining
    cost $h^{*}$.}\label{tab:story-h}
  \begin{tabular}{@{}lcccccccc@{}}
    \toprule
    City & \cty{D} & \cty{G} & \cty{N} & \cty{K} & \cty{B} & \cty{C} & \cty{H} & \cty{S}\\
    \midrule
    Heuristic $h$        & 7 & 8 & 5 & 4 & 4 & 7 & 2 & 0\\
    True cost $h^{*}$    & 8 & 9 & 6 & 7 & 5 & 8 & 3 & 0\\
    \bottomrule
  \end{tabular}
\end{table}
